% Source presentations for register row N-c-005.

\subsection{The classical decision problem in B\"orger--Gr\"adel--Gurevich}
\label{sec:p09-presentation-n-c-005}

The source classifies prenex first-order satisfiability only after fixing the
vocabulary.  If \(\Pi\) is a set of quantifier prefixes, the notation
\([\Pi,p,f]\) denotes the equality-free prenex sentences whose prefixes lie
in \(\Pi\), with the numbers of predicate and function symbols of each arity
bounded by the sequences \(p\) and \(f\).  A superscript \({}^{=}\) admits
logical equality.  Constants and free variables are not silently included in
the base notation; they enter only through explicit translations
\cite[def.~1.3.1]{BorgerGradelGurevich1997}.  The brackets therefore name a
syntactic class of sentences, not a logic with an unspecified vocabulary.

The book also defines a \emph{reduction class}: a syntactic class to which
arbitrary first-order sentences can be mapped effectively while preserving
satisfiability \cite[def.~2.1.6]{BorgerGradelGurevich1997}.  This is stronger
than saying that one neighbouring prefix is undecidable.  It explains why the
classification records both decidable islands and nearby classes already able
to encode the unrestricted decision problem.

\subsubsection{Bernays--Sch\"onfinkel--Ramsey}

The relational \(\exists^*\forall^*\) class is treated both with and without
equality.  If
\[
 \exists x_1\cdots\exists x_p
 \forall y_1\cdots\forall y_m\,\theta
\]
is satisfiable and \(\theta\) is quantifier-free over a relational vocabulary,
choose witnesses for the existential block and take the induced substructure
on them; if \(p=0\), retain one arbitrary element.  A universal relational
matrix is preserved by induced substructures, so the resulting model has size
at most \(\max(1,p)\), even when equality is admitted
\cite[prop.~6.2.17]{BorgerGradelGurevich1997}.  This complete small-model proof
gives decidability and, with the source's matching upper- and lower-bound
arguments, NEXPTIME-completeness \cite[pp.~257--261]{BorgerGradelGurevich1997}.
Positive-arity function symbols would invalidate the induced-substructure
step; leading existential variables may instead be represented by constants
only through the stated equisatisfiable translation.

\subsubsection{G\"odel--Kalm\'ar--Sch\"utte}

The equality-free relational core is \(\forall^2\exists^*\), extended by a
leading existential block to \(\exists^*\forall^2\exists^*\).  Theorem~6.2.23
proves the finite-model property for the core, and Corollary~6.2.38 transfers it
to the leading-existential closure.  The source uses a probabilistic finite-
model construction, explicitly replacing older and more elaborate
constructions; the argument is a source-specific proof route, not the BSR
witness-substructure argument \cite[pp.~261--269]{BorgerGradelGurevich1997}.
The two adjacent universal variables and the absence of equality are essential.
With equality, the corresponding G\"odel prefix crosses the Goldfarb
undecidability boundary recorded in Table~\ref{tab:p09-prefix-classes}.

\subsubsection{Ackermann and the maximal boundaries}

The relational Ackermann prefix is \(\exists^*\forall\exists^*\).  The book
first gives a deterministic procedure for the monadic case and then proves a
general upper bound of
\[
 2^{c(n/\log n)^2}
\]
for a constant \(c\) in its formula-size measure, followed by lower-bound
analysis \cite[thms.~6.3.24 and 6.3.26]{BorgerGradelGurevich1997}.  This
algorithm is not obtained by treating the one universal quantifier as a
two-universal GKS prefix.

The historical synthesis on pp.~310--311 gives two different maximality
statements.  Without equality, Bernays--Sch\"onfinkel--Ramsey and
G\"odel--Kalm\'ar--Sch\"utte are the two maximal decidable prefix classes;
with equality, Bernays--Sch\"onfinkel--Ramsey and Ackermann are the two maximal
decidable classes.  Already without equality, \(\forall^3\exists\) and
\(\forall\exists\forall\) are reduction classes
\cite[pp.~310--311]{BorgerGradelGurevich1997}.  A prefix table which omits its
equality column therefore changes the mathematical classification.

Skolemization does not preserve membership in these function-free input
classes.  Trailing existentials in the Ackermann row generally produce unary
Skolem functions, and in the GKS row can produce unary or binary Skolem
functions according to the preceding universal variables.  Those Skolem forms
may be equisatisfiable, but they are outside the classified relational class.
Likewise, passing from validity to satisfiability negates the matrix and dualizes
each quantifier in place; it does not reverse the prefix order.

\paragraph{Lewis 1979 source boundary.}
The register separately names Lewis's \emph{Unsolvable Classes of
Quantificational Formulas}.  The book's identity is established, but its
technical chapters were not available for verification.  Its exact class
notation, equality/function conventions, reductions, proofs, and examples are
therefore not supplied by copying this later 1997 synthesis.  Lewis's distinct
1980 complexity article is also not a substitute for the registered 1979 book
\cite{Lewis1979,Lewis1980}.
